\documentclass[11pt]{article}
\usepackage[margin=1in]{geometry}
\usepackage{amsmath,amssymb}
\usepackage[hidelinks]{hyperref}
\usepackage[round]{natbib}

\title{Asymmetric Capital Control: Short Summary}
\date{}

\begin{document}
\maketitle
\thispagestyle{empty}

\noindent
Codes in parentheses refer to results in the proofs document, which
carries an index of them.

\bigskip

Two banks facing identical regulation can run visibly different capital
policies. A bank that treats a shortfall as worse than an equivalent
\textbf{surplus} is good will hold more capital, pay out later, and
recapitalise sooner than the rules alone require---but so will a bank
simply responding to tight regulation. The two are not distinguishable by
inspection, yet they predict different behaviour when the rules change.

This paper solves a model in which they can be told apart, and the natural
place to look is the wrong one. A bank manages its capital \textbf{buffer}
under a \textbf{regulatory cap}, paying out at a \textbf{dividend barrier} and
raising capital in lump sums at a \textbf{recapitalisation trigger}, with
shortfalls below a \textbf{reference level} penalised by one asymmetry
parameter. At its neutral value the problem exactly nests the
\textbf{risk-neutral benchmark}~(NST), so the parameter captures asymmetry
beyond what regulation and costs alone produce.

The \textbf{definitional evaluation} of the \textbf{deficit}-to-\textbf{surplus} volatility
ratio---the second-moment measure one would instinctively reach for---is a
constant of \textbf{cap geometry}, carrying no information about the preference
whatsoever~(SCG). The thresholds carry it, and both responses are obtained
in closed form: a more shortfall-averse bank holds a larger \textbf{buffer}
before paying out and recapitalises earlier~(TCS). The \textbf{fixed issuance
cost} moves the \textbf{recapitalisation trigger} the opposite way~(KCS), and
that opposition is decisive --- it makes the Jacobian of the two thresholds
in the two parameters non-singular, so preference and cost are jointly
identified from where the bank commits to act~(IDN). The identification turns on a sign, not on a
comparison of magnitudes.

Preference is legible in the boundaries a bank commits to, not in the
volatility its behaviour produces. To our knowledge no firm-level model
combines singular payout control with impulse \textbf{recapitalisation} under
an asymmetric preference; the recent preference-generalised dividend
literature remains singular-only \citep{chen2026,cao2026}.

\bibliographystyle{plainnat}
\bibliography{references}

\end{document}
